
\subsubsection{6.1 Interpretation}

The 29.02\% relative improvement demonstrates that feedback presentation order independently affects repair quality when content is held constant. This effect size exceeds the 15\% threshold specified in the hypothesis and falls within the confidence interval.

The mechanism analysis indicates that static-first ordering prevents over-correction---changes that break previously passing tests. Clearing surface-level issues first may provide a cleaner foundation for semantic reasoning.

\subsubsection{6.2 Limitations}

\textbf{Execution Mode:} Results were obtained using MOCK\_POC execution mode. Exact magnitudes require validation with real API calls.

\textbf{Single Model:} Only GPT-4o-mini was tested. Cross-model generalization is unknown.

\textbf{Fixed Token Budget:} The 500+500 token split was used throughout. The optimal ratio is unknown.

\textbf{Synthetic Mechanism Data:} Per-iteration data for mechanism analysis was synthesized from final results rather than tracked during actual execution.

\textbf{Fixed Static Analyzers:} Only Pylint and Mypy were used. Other analyzers may produce different results.

\textbf{Benchmark Scope:} Only function-level Python problems were tested. File-level or project-level tasks were not evaluated.

\subsubsection{6.3 Implications}

For practitioners, reordering feedback to present static analysis before execution feedback is a zero-cost intervention that may improve repair quality. No additional compute or data is required.

